\section{Introduction}
\label{sec:introduction}

Fire safety engineers describe the growth of a fire with a design fire: a curve of heat release rate against time.
The most common form is the t-squared fire, in which the heat release rate grows with the square of the time since
ignition~\cite{heskestad1982}. Its growth coefficient is chosen from four classes, slow, medium, fast and ultrafast,
according to the fuel in the room.

A design fire drives every later calculation: the flame height, the temperature of the ceiling jet and the time at
which a sprinkler or a smoke detector responds~\cite{alpert1972}. Physics-based models such as the Wildland-urban
interface Fire Dynamics Simulator (WFDS)~\cite{mell2009}can represent the same processes in three dimensions, at a
far greater cost.

This note follows a house fire in the Castellamare neighborhood, where the fire spread structure-to-struture before
the first engine arrived. Section~\ref{sec:design-fires} describes the design fires, and the results follow in the
next section; the limits of the method are discussed in Section 4.0.
